\begin{definition}[The BRB implementation instance]
  \label{def:aba-bracha-implementation}
  \lean{PLTS.ABA.BRB.Message}\lean{PLTS.ABA.BRB.ProcessVariables}\lean{PLTS.ABA.BRB.BrachaState}\lean{PLTS.ABA.BRB.BroadcastEvent}\lean{PLTS.ABA.BRB.brachaInstanceExtended}\lean{PLTS.ABA.BRB.BrachaAlgorithm}\lean{PLTS.ABA.BRB.brachaInstance}\lean{PLTS.ABA.BRB.brachaInstance_step_iff_algorithm}\leanok
  \uses{def:aba-brb-spec, def:aba-networkstate, def:synchronised-product, def:process-algebra, def:relabel}
  Bracha's algorithm~\cite{Bra87} --- Algorithm~1 of~\cite{AFW25} (D34) --- over the two-part state of Definition~\ref{def:aba-networkstate}.
  The instance is the $n$ per-process programs under $\mathrm{synchronisedProduct}$ beside the instance's network, with the multicast and the delivery hidden: \leandecl{PLTS.ABA.BRB.brachaInstance}.
  A program holds one process's variables and the messages delivered to it, indexed by sender; it holds no sent set and no corrupted set, and its guards read its own variables and its own delivered sets.
  The network holds the per-sender sent sets and the corrupted set and reads no program's variables.
  A multicast is a synchronised step of the sender, which writes its variables, and the network, which records the message; a delivery is a synchronised step of the network, which checks that the message is sent under the named sender, and the receiver, which files it under that sender.
  The message type is $\mathrm{Message}(M) = \{\langle\mathrm{INIT}, m\rangle, \langle\mathrm{ECHO}, m\rangle, \langle\mathrm{VOTE}, m\rangle : m \in M\}$, and a process's variables hold its input beside the write-once send fields $\mathit{sentInit}$, $\mathit{sentEcho}$ and $\mathit{sentVote}$ and its return flag.
  What leaves the instance is the call at the leader, the per-process returns and corruption, the alphabet of Definition~\ref{def:aba-brb-spec} (D5, D1).
  The call writes the leader's input and nothing else, and the network records nothing at it.
  The specification takes its call label on two transitions, the call and the input-enabledness loop, and the leader's program carries the same two: \leandeclnamed{call}{PLTS.ABA.BRB.ProgramStep.call} records the payload at a leader holding no input, and \leandeclnamed{callLoop}{PLTS.ABA.BRB.ProgramStep.callLoop} writes nothing at a leader holding one.
  Every other program and the network take the call on a Dirac self-loop.
  \leandecl{PLTS.ABA.BRB.brachaInstance_step_iff_algorithm} is the characterisation: at every label, the transitions of the composition on that label are exactly the transitions of \leandecl{PLTS.ABA.BRB.BrachaAlgorithm} on it, on the same state and with the same distribution.
  Those transitions are the levels of the algorithm.
  $\mathrm{INIT}\,m$ is a multicast of the leader, once, on its input $m$: the first line of the source, $\text{send}\ \langle\mathrm{init}, m\rangle\ \text{to all}$.
  $\mathrm{ECHO}\,m$ fires on a delivered $\langle\mathrm{INIT}, m\rangle$ from the leader, on a quorum of received $\mathrm{ECHO}\,m$ messages from more than $(n+f)/2$ senders, or on $f{+}1$ received $\mathrm{VOTE}\,m$ messages.
  $\mathrm{VOTE}\,m$ fires on that quorum of received $\mathrm{ECHO}\,m$ messages or on $f{+}1$ received $\mathrm{VOTE}\,m$ messages (the amplification transition), and the return of $m$ takes $2f{+}1$ received $\mathrm{VOTE}\,m$ messages, at a process that has not returned.
  The call and its loop are the two transitions of the algorithm on the call label.
\end{definition}
